\begin{frame}
    \frametitle{Научная новизна}
    \begin{itemize}
        \item Установлено, что взаимное влияние одномодового и межмодового четырёхволнового смешения в маломодовом оптическом волокне обусловлено наличием общей антистоксовой компоненты. Предложен метод подавления нежелательных спектральных компонент с помощью модового фильтра.
        
        \item Обнаружено, что в процессе формирования решётки квадратичной нелинейной восприимчивости при фотоиндуцированной генерации второй гармоники в германосиликатном волокне ширина кривых фазового синхронизма монотонно уменьшается, их пик сдвигается к расстройке, определяемой фазовой самомодуляцией и кросс-модуляцией, при этом форма кривых синхронизма сохраняет выраженную асимметрию.  
        
        \item Для экспоненциального, импульсного и градиентного методов обработки температурных кинетик предложены поправочные выражения, позволяющие компенсировать нестационарность температуры окружающей среды и коэффициента теплообмена.
    \end{itemize}
\end{frame}
\note{
    Проговаривается вслух научная новизна
}

\begin{frame}
    \frametitle{Практическая значимость}
    \begin{itemize}
        \item Предложенный метод подавления межмодового четырёхволнового смешения с помощью модового фильтра позволяет увеличить допустимую длину волокна доставки излучения на 0.5~м без снижения эффективности генерации разностной частоты, что важно для создания эффективных и надёжных гибридных источников среднего ИК-диапазона.
        
        \item Полученные зависимости эволюции кривых фазового синхронизма при фотоиндуцированной ГВГ позволяют понимать, как меняется фазовая расстройка в процессе формирования решётки. Это создаёт основу для будущих исследований по управлению фазой накачки с целью либо подавления паразитной ГВГ, либо её эффективной генерации.
        
        \item Разработанные поправочные выражения для обработки температурных кинетик позволяют снизить разброс измерений коэффициента оптического поглощения в три раза, повышая точность диагностики нелинейно-оптических кристаллов и достоверность оценки их теплового режима при работе в мощных лазерных системах
    \end{itemize}
\end{frame}
\note{
    Проговариваются вслух научная и практическая значимость
}


\begin{frame} % публикации на одной странице
% \begin{frame}[t,allowframebreaks] % публикации на нескольких страницах
    \frametitle{Основные публикации}
    \nocite{vakbib1}%
    \nocite{vakbib2}%
    %
    %% authorwos
    \nocite{wosbib1}%
    %
    %% authorscopus
    \nocite{scbib1}%
    %
    %% authorconf
    \nocite{confbib1}%
    \nocite{confbib2}%
    %
    %% authorother
    \nocite{bib1}%
    \nocite{bib2}%
    \ifnumequal{\value{bibliosel}}{0}{
        \insertbiblioauthor
    }{
        \printbibliography%
    }
\end{frame}
\note{
    Результаты работы опубликованы в N печатных изданиях,
    в~т.\:ч. M реферируемых изданиях.
}

\begin{frame}
    \frametitle{Участие в конференциях}
    \begin{itemize}
		\item 65-я, 66-я, 67-я Всероссийская научная конференция МФТИ (Москва 2023, 2024, 2025)
		\item International Conference on Advanced Laser Technologies (ALT) (Казань 2025)
		\item XX Международная молодежная конференция по люминесценции и лазерной физике (LLPh) (Иркутск 2023, 2025)
		\item International Conference Laser Optics (ICLO) (Санкт-Петербург 2022, 2024)
		\item XII Международная конференция по фотонике и информационной оптике (Москва 2023)
		\item XXI Всероссийская молодежная Самарская конкурс-конференция по оптике, лазерной физике и физике плазмы, посвященная 300-летию РАН (Москва 2023)
		\item XVII Молодежный конкурс имени Ивана Анисимкина (Москва 2021)
		
    \end{itemize}
\end{frame}
\note{
    Работа была представлена на ряде конференций.
}

\begin{frame}[plain, noframenumbering] % последний слайд без оформления
    \begin{center}
        \Huge
        Спасибо за внимание!
    \end{center}
\end{frame}
